\documentclass[../../main.tex]{subfiles}% 注意这里的文件路径不能用 ./main.tex ,否则用latexmk编译子文件会报错
\graphicspath{{\subfix{./image/}}{\subfix{../../image/}}} % 指定图片引用路径,后续可以直接使用图片文件名
% 如果图片存放在上级目录,则使用 ../../image/ 作为备用路径

% 例如:
% \begin{figure}[H]
% \centering
% \includegraphics[width=0.3\textwidth]{图.png}
% \caption{}
% \label{figure:图}
% \end{figure}
% 注意:上述\label{}一定要放在\caption{}之后,否则引用图片序号会只会显示??.

\begin{document}

\section{唯一性}\label{section:ch06-唯一性}

\cref{theorem:ch06-6.2.1}中的持续性性质在证明 (P) 弱解唯一性时起重要作用，因为它完全刻画了弱解的 $2$-集中部分。另一方面，由于\cref{definition:ch06-弱解}对 $2$-弥散部分 $[u_s(\cdot,t)]_{d,2}$ 未作任何规定（见\eqref{align:ch03-3.82}），不出所料，(P) 的弱解不构成问题的唯一类（见\cref{remark:ch06-6.3.2}）。这使如下定义变得重要。

\begin{definition}\label{definition:ch06-相容性条件}
设 (A0)--(A1) 成立。(P) 的弱解 $u$ 满足相容性条件，若对几乎所有 $t\in(0,T)$，
\begin{align}
[u^\pm_s(\cdot,t)]_{d,2}=[u^\pm_s(\cdot,t)]_{d,2}\llcorner\mathscr{S}^\pm_t, \label{align:ch6-6.21a}
\end{align}
其中
\begin{align}
\mathscr{S}^\pm_t\triangleq \{x\in\Omega\mid\phi^\pm(u_r(x,t))=\pm\phi_{\pm\infty}\}. \label{align:ch6-6.21b}
\end{align}
\end{definition}

在本定义及此后，凡符号 $\pm$ 与 $\mp$ 出现的陈述，均应读作一对独立陈述，分别对应上、下号选择。

\begin{remark}\label{remark:ch06-6.3.1}
由\cref{definition:ch06-弱解}(i)，对几乎所有 $t\in(0,T)$，$\phi(u_r(\cdot,t))\in H^1_0(\Omega)$，故 $\phi(u_r(\cdot,t))$ 有一个 $2$-拟连续代表元 $v(\cdot,t)\in H^1_0(\Omega)$（见\cref{definition:ch02-拟连续}）。在\eqref{align:ch6-6.21b}中，我们将函数 $\phi^\pm(u_r(\cdot,t))$ 等同于其 $2$-拟连续代表元 $v^\pm(\cdot,t)\in H^1_0(\Omega)$，后者在 $\mathbb{R}$ 上 $2$-拟处处满足 $0\leqslant v^\pm\leqslant\pm\phi_{\pm\infty}$。在\eqref{align:ch6-6.21a}中，我们总将集 $\mathscr{S}^\pm_t$ 等同于 $\widetilde{\mathscr{S}}^\pm_t\triangleq \{x\in\Omega\mid v^\pm(x,t)=\pm\phi_{\pm\infty}\}$。因 $v^\pm(\cdot,t)$ 定义到相差一个 $2$-零容集，集 $\widetilde{\mathscr{S}}^\pm_t$ 亦然。
\end{remark}

问题 (P) 的唯一性由如下定理给出，它利用\cref{definition:ch06-相容性条件}中的相容性条件。

\begin{theorem}\label{theorem:ch06-6.3.1}
设 (A0)--(A1) 成立。则对每个 $u_0\in\mathscr{R}_f(\Omega)$，问题 (P) 至多存在一个满足相容性条件\eqref{align:ch6-6.21a}的弱解。
\end{theorem}

\begin{remark}\label{remark:ch06-6.3.2}
满足相容性条件的 (P) 弱解的存在性将在下文证明（见\cref{theorem:ch06-6.4.2}）。另一方面，容易构造不满足相容性条件的 (P) 弱解，因而它们不可能是\cref{theorem:ch06-6.4.2}所给出的解。故弱解不构成问题 (P) 的唯一类。

设 (A0) 成立且 $u_{0s}\neq0$。由
\begin{align}
u(\cdot,t)\triangleq w(\cdot,t)+u_{0s} \quad\text{对几乎所有 }t\in(0,T) \label{align:ch6-6.22}
\end{align}
定义 $u\in L^\infty_{w^*}(0,T;\mathscr{R}_f(\Omega))$，其中 $w\in L^\infty(0,T;L^1(\Omega))$ 为初值 $u_{0r}$ 对应的 (P) 唯一弱解（$w$ 的存在唯一性由\cref{theorem:ch06-6.4.3}保证）。显然 $u$ 是 (P) 的弱解，$u_r=w$ 在 $Q$ 中几乎处处成立，奇异部分 $u_s(\cdot,t)=u_{0s}$ 在时间上为常数，因而 $u$ 满足相容性条件当且仅当
\begin{align*}
[u^\pm_{0s}]_{d,2}=[u^\pm_{0s}]_{d,2}\llcorner\{x\in\Omega\mid\phi^\pm(w(x,t))=\pm\phi_{\pm\infty}\}
\quad\text{对几乎所有 }t\in(0,T).
\end{align*}
这些等式不成立的例子容易给出，故断言成立。
\end{remark}

\begin{proof}[证明\cref{theorem:ch06-6.3.1}]
设 $u_1$ 与 $u_2$ 为问题 (P) 的两个弱解。为简便记 $u_{1r}\equiv(u_1)_r$，$u_{1s}\equiv(u_1)_s$（$u_2$ 类似）。设 $u_1$ 与 $u_2$ 满足相容性条件
\begin{align}
[u^\pm_{is}(\cdot,t)]_{d,2}=[u^\pm_{is}(\cdot,t)]_{d,2}\llcorner\mathscr{S}^t_{i,\pm}, \label{align:ch6-6.23}
\end{align}
其中
\begin{align*}
\mathscr{S}^t_{i,\pm}\triangleq \{x\in\Omega\mid\phi(u_{ir})(x,t)=\phi_{\pm\infty}\}\quad(i=1,2).
\end{align*}
由\eqref{align:ch6-6.2}立即得
\begin{align}
\int_0^T\langle u_1(\cdot,t)-u_2(\cdot,t),\partial_t\zeta(\cdot,t)\rangle\,\mathrm{d}t
= \iint_Q\nabla[\phi(u_{1r})-\phi(u_{2r})]\cdot\nabla\zeta\,\mathrm{d}x\,\mathrm{d}t \label{align:ch6-6.24}
\end{align}
对每个 $\zeta\in C^1([0,T];C^1_c(\Omega))$，$\zeta(\cdot,T)=0$ 在 $\Omega$ 上成立。又由等式\eqref{align:ch6-6.12}，对所有 $t\in(0,T)$，有
\begin{align}
[u_1(\cdot,t)]_{c,2}=[u_0]_{c,2}=[u_2(\cdot,t)]_{c,2} \quad\text{在 }\mathscr{R}_f(\Omega). \label{align:ch6-6.25}
\end{align}
由\eqref{align:ch6-6.24}--\eqref{align:ch6-6.25}得
\begin{align}
\int_0^T\langle[u_1(\cdot,t)]_{d,2}-[u_2(\cdot,t)]_{d,2},\partial_t\zeta(\cdot,t)\rangle\,\mathrm{d}t
= \iint_Q\nabla[\phi(u_{1r})-\phi(u_{2r})]\cdot\nabla\zeta\,\mathrm{d}x\,\mathrm{d}t \label{align:ch6-6.26}
\end{align}
对每个 $\zeta\in C^1([0,T];C^1_c(\Omega))$，$\zeta(\cdot,T)=0$ 在 $\Omega$ 上成立。

由标准正则化论证，此等式对任意 $\zeta\in C([0,T];H^1_0(\Omega))$，$\partial_t\zeta\in L^\infty(Q)\cap L^2(0,T;H^1_0(\Omega))$，$\zeta(\cdot,T)=0$ 几乎处处在 $\Omega$ 上也成立。这使我们可在\eqref{align:ch6-6.26}中选取试验函数
\begin{align*}
\zeta(x,t)\triangleq -\int_t^T[\phi(u_{1r})-\phi(u_{2r})](x,s)\,\mathrm{d}s \quad((x,t)\in Q),
\end{align*}
从而得到
\begin{align}
\begin{aligned}
&\int_0^T\langle[u_1(\cdot,t)]_{d,2}-[u_2(\cdot,t)]_{d,2},
[\phi(u_{1r})-\phi(u_{2r})](\cdot,t)\rangle\,\mathrm{d}t\\
&\quad= -\iint_Q\nabla[\phi(u_{1r})-\phi(u_{2r})]\cdot
\left(\int_t^T\nabla[\phi(u_{1r})-\phi(u_{2r})](x,s)\,\mathrm{d}s\right)\mathrm{d}x\,\mathrm{d}t.
\end{aligned} \label{align:ch6-6.27}
\end{align}
由等式\eqref{align:ch03-3.82}有
\begin{align}
\begin{aligned}
&\int_0^T\langle[u_1(\cdot,t)]_{d,2}-[u_2(\cdot,t)]_{d,2},
[\phi(u_{1r})-\phi(u_{2r})](\cdot,t)\rangle\,\mathrm{d}t\\
&\quad= \int_0^T\langle[u_{1s}(\cdot,t)]_{d,2}-[u_2s(\cdot,t)]_{d,2},
[\phi(u_{1r})-\phi(u_{2r})](\cdot,t)\rangle\,\mathrm{d}t\\
&\qquad+\iint_Q(u_{1r}-u_{2r})[\phi(u_{1r})-\phi(u_{2r})]\,\mathrm{d}x\,\mathrm{d}t,
\end{aligned} \label{align:ch6-6.28}
\end{align}
而由相容性条件\eqref{align:ch6-6.23}，
\begin{align}
\begin{aligned}
&\int_0^T\langle[u_{1s}(\cdot,t)]_{d,2}-[u_{2s}(\cdot,t)]_{d,2},
[\phi(u_{1r})-\phi(u_{2r})](\cdot,t)\rangle\,\mathrm{d}t\\
&\quad= \int_0^T\langle[u^+_{1s}(\cdot,t)]_{d,2},\phi_{+\infty}-\phi(u_{2r})(\cdot,t)\rangle\,\mathrm{d}t\\
&\qquad-\int_0^T\langle[u^-_{1s}(\cdot,t)]_{d,2},\phi_{-\infty}-\phi(u_{2r})(\cdot,t)\rangle\,\mathrm{d}t\\
&\qquad-\int_0^T\langle[u^+_{2s}(\cdot,t)]_{d,2},\phi(u_{1r})(\cdot,t)-\phi_{+\infty}\rangle\,\mathrm{d}t\\
&\qquad+\int_0^T\langle[u^-_{2s}(\cdot,t)]_{d,2},\phi(u_{1r})(\cdot,t)-\phi_{-\infty}\rangle\,\mathrm{d}t
\geqslant0.
\end{aligned} \label{align:ch6-6.29}
\end{align}
由\eqref{align:ch6-6.28}--\eqref{align:ch6-6.29}得
\begin{align}
\begin{aligned}
&\int_0^T\langle[u_1(\cdot,t)]_{d,2}-[u_2(\cdot,t)]_{d,2},
[\phi(u_{1r})-\phi(u_{2r})](\cdot,t)\rangle\,\mathrm{d}t\\
&\quad\geqslant\iint_Q(u_{1r}-u_{2r})[\phi(u_{1r})-\phi(u_{2r})]\,\mathrm{d}x\,\mathrm{d}t.
\end{aligned} \label{align:ch6-6.30}
\end{align}
关于\eqref{align:ch6-6.27}右端，有
\begin{align}
\begin{aligned}
&-\iint_Q\nabla[\phi(u_{1r})-\phi(u_{2r})]\cdot
\left(\int_t^T\nabla[\phi(u_{1r})-\phi(u_{2r})](x,s)\,\mathrm{d}s\right)\mathrm{d}x\,\mathrm{d}t\\
&\quad= \frac{1}{2}\iint_Q\frac{\mathrm{d}}{\mathrm{d}t}
\left|\int_t^T\nabla[\phi(u_{1r})-\phi(u_{2r})](x,s)\,\mathrm{d}s\right|^2\mathrm{d}x\,\mathrm{d}t\\
&\quad= -\frac{1}{2}\int_\Omega\left|\int_0^T\nabla[\phi(u_{1r})-\phi(u_{2r})](x,s)\,\mathrm{d}s\right|^2\mathrm{d}x.
\end{aligned} \label{align:ch6-6.31}
\end{align}
由\eqref{align:ch6-6.27}与\eqref{align:ch6-6.30}--\eqref{align:ch6-6.31}得
\begin{align*}
\iint_Q(u_{1r}-u_{2r})[\phi(u_{1r})-\phi(u_{2r})]\,\mathrm{d}x\,\mathrm{d}t
\leqslant-\frac{1}{2}\int_\Omega\left|\int_0^T\nabla[\phi(u_{1r})-\phi(u_{2r})](x,t)\,\mathrm{d}t\right|^2\mathrm{d}x\leqslant0,
\end{align*}
从而
\begin{align}
\phi(u_{1r}(x,t))=\phi(u_{2r}(x,t)) \quad\text{对几乎所有 }(x,t)\in Q. \label{align:ch6-6.32}
\end{align}
现在在等式\eqref{align:ch6-6.24}中取 $\zeta(x,t)=\rho(x)k_q(t)$，其中 $\rho\in C^2_c(\Omega)$，$k_q$ 如\eqref{align:ch6-6.7}。如\eqref{align:ch6-6.4}证明中那样论证，并利用等式\eqref{align:ch6-6.32}，对几乎所有 $t_1\in(0,T)$，得
\begin{align*}
\langle u_1(\cdot,t_1)-u_2(\cdot,t_1),\rho\rangle=0.
\end{align*}
由 $\rho$ 的任意性，这蕴含对几乎所有 $t_1\in(0,T)$，$u_1(\cdot,t_1)=u_2(\cdot,t_1)$ 在 $\mathscr{R}_f(\Omega)$。结论随之成立。
\end{proof}




\end{document}